\textbf{Цель работы}
\begin{itemize}
	\item познакомиться с методологией CI/CD; 
	\item научиться работать с инструментом GitHub Actions. 
\end{itemize}

\textbf{Задачи:}
\begin{enumerate}
	\item изучить методические материалы к лабораторной работе; 
	\item написать тесты для приложения на Python; 
	\item написать пайплайн сборки, тестирования и загрузки образа приложения на Docker Hub. 
\end{enumerate}

\section*{Выполнение задания}

%cd devops-lab5
%git remote set-url origin https://github.com/yeeltheeel/devops-mtuci-lab5
%git checkout -b dev
%mkdir .github/workflows
%cp ../tests.yml ../build-and-delivery.yml .github/workflows
%git add .github/workflows
%git commit -m "added files"
%git push --set-upstream origin dev

\begin{figure}[!htbp]
	\centering
	\includegraphics[width=0.6\linewidth]{"Screenshot 2026-04-09 021110"}
	\caption{Новый репозиторий с веткой dev}
\end{figure}
\FloatBarrier

\begin{figure}[!htbp]
	\centering
	\includegraphics[width=0.6\linewidth]{"Screenshot 2026-04-08 222438"}
	\caption{Работа пайплайна в ветке dev}
\end{figure}
\FloatBarrier

\begin{figure}[!htbp]
	\centering
	\includegraphics[width=0.6\linewidth]{"Screenshot 2026-04-09 020939"}
	\caption{Работа пайплайнов в ветке main}
\end{figure}
\FloatBarrier

\begin{figure}[!htbp]
	\centering
	\includegraphics[width=0.6\linewidth]{"Screenshot 2026-04-09 022419"}
	\caption{Загруженный образ в Docker Hub}
\end{figure}
\FloatBarrier

%docker pull ghcr.io/yeeltheeel/my-app:latest
%docker run -p 8080:8080 ghcr.io/yeeltheeel/my-app
%http://localhost:8080/docs
\begin{figure}[!htbp]
	\centering
	\includegraphics[width=0.6\linewidth]{"Screenshot 2026-04-09 022218"}
	\caption{Запущенный образ}
\end{figure}
\FloatBarrier

\section*{Исходный код программы}

%\lstinputlisting[style=py]{devops-lab5/tests/test_user.py}
\lstinputlisting[style=py, linerange={1, 3-6, 8-21, 23-27, 29-32, 34-41, 43-50, 52-58}]{devops-lab5/tests/test_user.py}
Листинг 1 -- Дополненные методы тестирования   

\lstinputlisting[style=ymlcicd]{devops-lab5/.github/workflows/tests.yml}
Листинг 2 -- Пайплайн тестирования

%\lstinputlisting[style=ymlcicd]{devops-lab5/.github/workflows/build-and-delivery.yml}
\lstinputlisting[style=ymlcicd, linerange={1-11, 16-18, 21-27, 29-32}]{devops-lab5/.github/workflows/build-and-delivery.yml}
Листинг 3 -- Пайплайн сборки и загрузки образа приложения
%ghcr.io -- github container registry (store/pull)

\section*{Заключение}

В ходе выполнения лабораторной работы были решены следующие задачи:
\begin{itemize}
	\item проведено ознакомление с методологией CI/CD; 
	\item получены знания о том, как работать с инструментом GitHub Actions. 
\end{itemize}

\section*{Контрольные вопросы}

\begin{enumerate}
	\item \textit{Как CI/CD помогает в разработке?} Помогает избежать ошибок и сбоев в коде, поддерживая непрерывный цикл разработки и обновления программного обеспечения. Поскольку CI/CD автоматизирует ручное вмешательство человека, время простоя сводится к минимуму, а релизы кода происходят быстрее, снижается человеческий фактор.
	\item \textit{К какому этапу CI/CD относится сборка десктоп приложения под операционные системы Windows, MacOS и Linux для его дальнейшего тестирования?} CI.
	\item \textit{Что такое Continuous Delivery и Deploy?} Continuous Delivery (непрерывная доставка) автоматизирует выпуск проверенного кода в репозиторий после сборки и тестирования приложения на этапе CI. Deploy (непрерывное развертывание) -- заключительный этап конвейера CI/CD, является расширением непрерывной доставки и может означать автоматизацию выпуска изменений из репозитория в производство.
\end{enumerate}
